% Authored reconciliation of existing results, 2026-11-05. No ledger promotion.
\section{Shared origin of the seam normalization and carrier data}
\label{sec:shared-p1-p2-origin}

\noindent\textbf{What is being explained.} P1 supplies a primitive reflection-positive
kernel on an oriented seam with unit winding. The number $c_3$ is its normalization.
P2 declares the five-slot $3+2$ readout; its algebraic consequences and its
conditional source reductions are already available. Calling both entries
``axioms'' records these interfaces and does not erase those reductions.

\paragraph{One marked geometry, several typed readouts.}
In the flat spherical $\mathbb Z_2$ branch class, each branch point has deficit
$\pi$. Gauss--Bonnet gives $n\pi=4\pi$, hence four marks
(\veri{v216\_marks\_gauss\_bonnet.py}). For the common divisor
$D=\mu_4\subset\mathbb P^1$,
\[
 \operatorname{rank}H_1(\mathbb P^1\setminus D)=3,
 \qquad h^0(\mathcal O(D))=5.
\]
The functions and cycles belong to the same marked geometry
(\veri{v189\_riemann\_roch\_carrier.py},
\veri{v197\_rr\_carrier\_clifford\_d5.py}). They are different spaces:
meromorphic functions have four independent principal parts, while
meromorphic differentials have the residue-sum constraint and dimension three.
These dimensions alone do not identify a physical colour/weak $3+2$ split.
The square position of the marks is additional to their count; in the
flag-symmetric branch it is equivalent to $\tau=i$
(\veri{v512\_seam\_tau\_flag.py}).

\paragraph{The anchor and the two numerical readouts.}
The typed parabolic splitting $a=(1,1,2)$ gives
\[
 e_1(a)=4,\quad e_2(a)=5,\quad e_3(a)=2,
 \qquad c_3=\frac{1}{2\pi e_1(a)}=\frac1{8\pi},
 \qquad g_{\rm car}=e_2(a)=5.
\]
The uniqueness of three positive integer weights summing to four is exact;
the Deligne extension, Birkhoff--Grothendieck splitting and stability route
explain the typing for the specified connection
(\veri{v491\_p2\_partition\_corollary.py},
\veri{v499\_p2\_weights\_deligne\_bg.py}). Later analytic descent of
the actual $A_0^\star$ system gives finite $S_4$ monodromy, so a positive
invariant Hermitian form exists by finite-group averaging. Historical
analytic-monodromy and unitarizability obligations are therefore not a
current independent numerical gap. Identifying that mathematical module
and its frame with the actual physical source retains its stated scope.\footnote{TFPT Gesamtsynthese, 28 September 2026, the fourfold cover section; recovered source proof in \texttt{REC20261105C3}.}

\paragraph{Two normalization arguments with different hypotheses.}
Unit seam winding gives $\Delta=2\pi$ and the relative Chern normalization
$c_3=\Delta/(16\pi^2)=1/(8\pi)$. The one-sided normal-sphere argument
gives the same value.\footnote{Boundary-kernel source, primitive seam theorem; original source paths and spans are listed in \texttt{REC20261105C3}.} The index/KMS argument instead uses
$I=4$, $\beta_{\rm angle}=2\pi$ and
$c_3=1/(I\beta_{\rm angle})$. Its arithmetic is exact
(\veri{v813\_p1\_index\_kms.py}); the determinant-line response must
also be identified with the additive transitive four-channel response.
The charged MMST paper explicitly retains this additional operator
identification even after its net/KMS rigidity theorem. Equal channel
weights alone do not supply the variation formula.

\paragraph{Which former prerequisites have already been derived.}
The W-bridge anchor direction is derived by the winding/triple lock
(\veri{v1014\_bridge\_refinements.py}). Its democratic response is
derived for the finite factorized abelian collar model
(\veri{v1015\_axiom\_core\_closure.py}): $r=(1,1,1)$ at sizes $12,16$,
with character-dependent and deck-breaking controls. Both finite W-bridge
premises are closed; its finite axiom core is zero modulo the externalized
MMST identification. This does not remove the continuum determinant
identification in the P1 route. The earlier MMST route retains its conditional ledger status. The later
charged master theorem proves its implication under microscopic local-algebra
exhaustion, which remains unproved for the frozen collar. The unconditional
parent retains its ledger marker.

The 2 November R102 result also removes evenness, unimodularity, fixed
rank eight and minimality as inputs to a stronger selection theorem:
for the standard positive $D_{n+1}\oplus A_{n-1}$ pairing, an integral
full-rank extension containing a primitive $A$ discriminant projection
forces $n\mid4$, hence $n=4$ for $n\ge3$, and the whole completion is
$E_8$. The standard positive normalization and the actual one-sided
integral primitive insertion remain its source hypotheses. Full doubled
locality alone does not imply that one-sided hypothesis. The complementary
metric rigidity theorem assumes unimodularity, so the two results cannot
be used to prove each other's premises in a circle.

\paragraph{The common explanation and its boundary.}
The marked seam supplies the geometric setting; its cycles, sections,
parabolic anchor and primitive normalization explain the shared $4,3,5,2$
data in their specified branches. P1, its scalar normalization and P2
are three levels of this construction. The operational original seed
already contains the local net, state, reflection, time flow and collar
generator. Existing conditional OS and common-action reconstructions must
be reused with those data. Source identification means proving the named
maps for the same source; it does not mean constructing state and time
anew from the scalar $c_3$. The consolidated original-source audit is
\texttt{REC20261105C3}; it changes no axiom or T1--T8 ledger marker.
